%! AP Statistics — Unit 2, Lesson 2.1 — Lesson Plan
\documentclass[10pt]{article}
\usepackage{apstats-article}
\usepackage{apstats-boxes}

\newcommand{\UnitNumberName}{Unit 2: Exploring Two-Variable Data \quad}
\newcommand{\LessonNumberName}{Lesson 2.1: Introducing Statistics --- Are Variables Related?}

\begin{document}
\begin{center}
    {\Large\bfseries \CourseName: \SchoolYear\ (\MeetingLength) \\
    {\small\bfseries \UnitNumberName \LessonNumberName}}
\end{center}

\begin{tcolorbox}[colback=sky, colframe=navy, boxrule=0.9pt, arc=2mm,
    left=2mm, right=2mm, top=1.5mm, bottom=1.5mm]
    \textbf{Primary Objective:} Students will identify statistical questions that can be
    answered using two-variable data and recognize that apparent patterns or associations
    in data may be random or non-random --- making conclusions uncertain (Big Idea: VAR).
\end{tcolorbox}

\renewcommand{\arraystretch}{1.6}
\begin{skillbox}[Priority Ideas \& Skills]{goldbox}
\begin{minipage}[t]{0.48\linewidth}
    \textbf{AP Skill 1 --- Selecting Statistical Methods}
    \begin{itemize}
        \item Identify a statistical question that asks about a possible relationship between two variables (Skill 1.A).
        \item Recognize when an apparent association may be due to chance rather than a real relationship (VAR-1.D).
        \item Distinguish between a question about one variable (Unit 1) and a question about the relationship between two variables (Unit 2).
    \end{itemize}
\end{minipage}
\hfill
\begin{minipage}[t]{0.48\linewidth}
    \textbf{Key Understandings}
    \begin{itemize}
        \item Apparent patterns in data are not always real --- they may be random fluctuation.
        \item The statistical question determines how we analyze and display data.
        \item Even a real association does not prove causation.
        \item Unit 2 builds directly on Unit 1: now we compare two variables across the same individuals.
    \end{itemize}
\end{minipage}
\end{skillbox}

\begin{skillbox}[{Vocabulary, Concepts, \& Theorems}]{greenbox}
\begin{minipage}[t]{0.48\linewidth}
    \begin{itemize}
        \item \textbf{Two-variable data} --- a dataset in which two variables are recorded for each individual
        \item \textbf{Association} --- a relationship between two variables such that knowing one gives information about the other
        \item \textbf{Apparent association} --- a pattern visible in a sample that may or may not reflect a real relationship in the population
    \end{itemize}
\end{minipage}
\hfill
\begin{minipage}[t]{0.48\linewidth}
    \begin{itemize}
        \item \textbf{Causation} --- one variable directly produces a change in another; cannot be established from observational data alone
        \item \textbf{Confounding variable} --- a third variable related to both, making an apparent association misleading
        \item \textbf{Explanatory / Response variable} --- preview: the variable used to explain (explanatory) vs.\ the outcome of interest (response)
    \end{itemize}
\end{minipage}
\end{skillbox}

\begin{skillbox}[Activate Prior Knowledge \& Spiral Review]{sky}
\begin{minipage}[t]{0.48\linewidth}
    \textbf{Prior Knowledge Check (3--4 min)}
    \begin{itemize}
        \item Quick-fire recall: ``What are the two types of variables from Unit 1?''
          (categorical, quantitative)
        \item Transition question: \textit{``Unit 1 described one variable at a time.
          What would change if we wanted to ask whether two variables are related?''}
    \end{itemize}
\end{minipage}
\hfill
\begin{minipage}[t]{0.48\linewidth}
    \textbf{Connections Forward}
    \begin{itemize}
        \item Two categorical variables $\to$ Lessons 2.2--2.3 (two-way tables, bar charts)
        \item Two quantitative variables $\to$ Lessons 2.4--2.9 (scatterplots, regression)
        \item Establishing causation $\to$ Unit 3 (experimental design)
        \item Making inferences about associations $\to$ Units 6--9
    \end{itemize}
\end{minipage}
\end{skillbox}

\begin{skillbox}[Hook \& Lesson]{sky}
\begin{minipage}[t]{0.48\linewidth}
    \textbf{Hook (5--7 min)}\\
    Display two headlines:
    \begin{itemize}
        \item \textit{``Students who eat breakfast score higher on tests.''}
        \item \textit{``Shoe size and reading ability are strongly associated in children.''}
    \end{itemize}
    Ask: Are these relationships real? Could they be random? Could something else explain them?
    Transition: \textit{``These are the questions Unit 2 is designed to answer.''}
\end{minipage}
\hfill
\begin{minipage}[t]{0.48\linewidth}
    \textbf{Lesson Flow (15--18 min)}
    \begin{enumerate}
        \item Define two-variable data; show a small dataset with two variables per individual.
        \item Ask: ``Does eating breakfast cause higher scores, or is something else going on?'' Introduce apparent vs.\ real association.
        \item Shoe-size example: children who wear larger shoes are older --- age is the lurking variable; the association is real but not directly causal.
        \item Preview Unit 2 road map: the type of analysis depends on the variable types.
    \end{enumerate}
\end{minipage}
\end{skillbox}

\begin{skillbox}[Explicit Instruction \& Active Monitoring]{sky}
\begin{minipage}[t]{0.48\linewidth}
    \textbf{Direct Instruction (10 min)}
    \begin{itemize}
        \item Walk through 3 research questions. For each, ask:
          (a) What two variables are involved?
          (b) Is the apparent pattern likely real or possibly random?
          (c) If real, does it prove causation?
        \item Emphasize EK VAR-1.D.1: an apparent association may be random.
          A sample of $n = 10$ can show a pattern that disappears with $n = 1{,}000$.
    \end{itemize}
\end{minipage}
\hfill
\begin{minipage}[t]{0.48\linewidth}
    \textbf{Active Monitoring}
    \begin{itemize}
        \item Circulate as students categorize question types on their guided notes.
        \item Listen for: students treating any pattern as a real relationship;
          address with ``could this just be random variation in a small sample?''
        \item Cold-call on the hook scenarios before moving to group work.
    \end{itemize}
\end{minipage}
\end{skillbox}

\begin{skillbox}[Group Work \& Differentiation]{redbox}
\begin{minipage}[t]{0.48\linewidth}
    \textbf{Partner Activity (8--10 min)}\\
    Provide 4 research questions. Partners decide:
    \begin{enumerate}
        \item What two variables are being compared?
        \item Is this a question about a relationship (Unit 2) or about one variable (Unit 1)?
        \item Could an apparent association here be due to chance? What would convince you it is real?
        \item Can this study establish causation? Why or why not?
    \end{enumerate}
\end{minipage}
\hfill
\begin{minipage}[t]{0.48\linewidth}
    \textbf{Differentiation}
    \begin{itemize}
        \item \textit{Support}: Provide a sentence frame --- ``The two variables are \underline{\hspace{1cm}} and \underline{\hspace{1cm}}. The apparent association may be random because \underline{\hspace{1cm}}.''
        \item \textit{Extension}: Students write their own research question, identify a potential confounding variable, and explain how they would design a study to eliminate it.
        \item \textit{ELL}: Pair headline excerpts with a vocabulary anchor chart.
    \end{itemize}
\end{minipage}
\end{skillbox}

\begin{skillbox}[Individual Work \& Assessment]{redbox}
\begin{minipage}[t]{0.48\linewidth}
    \textbf{Exit Ticket (5 min)}\\
    Scenario: \textit{``A study of 30 high school students finds that students who drink
    more coffee get higher grades.''} Students independently write:
    \begin{enumerate}
        \item The two variables being studied
        \item Whether this association is likely real or possibly random (and why sample size matters)
        \item One possible confounding variable
        \item Whether this proves coffee causes higher grades
    \end{enumerate}
\end{minipage}
\hfill
\begin{minipage}[t]{0.48\linewidth}
    \textbf{AP-Style Check}\\
    \textit{``A study found a strong positive association between the number of fire trucks
    sent to a fire and the amount of property damage. Which of the following is the most
    plausible explanation?''}
    \begin{enumerate}[label=(\Alph*)]
        \item Fire trucks cause property damage
        \item Property damage attracts more fire trucks
        \item The size of the fire is a confounding variable
        \item The association is due to measurement error
    \end{enumerate}
    Use student responses to assess readiness for Lesson 2.2.
\end{minipage}
\end{skillbox}

\begin{skillbox}[Reinforcement \& Extension]{goldbox}
\begin{minipage}[t]{0.48\linewidth}
    \textbf{Homework / Practice}
    \begin{itemize}
        \item For 4 news blurbs reporting associations, identify the two variables, evaluate whether the association is likely real or possibly random, and identify one confounding variable.
        \item Reflect: \textit{``Why isn't a strong association between two variables enough to prove causation?''}
    \end{itemize}
\end{minipage}
\hfill
\begin{minipage}[t]{0.48\linewidth}
    \textbf{Extension / Preview}
    \begin{itemize}
        \item \textit{Extension}: Look up ``spurious correlations'' (tylervigen.com) and choose one. Explain why the association is almost certainly random, not real.
        \item \textit{Preview}: Lesson 2.2 begins analyzing relationships between \emph{two categorical} variables using a two-way frequency table --- start thinking about what such a table would look like.
    \end{itemize}
\end{minipage}
\end{skillbox}

\end{document}
